%% kapa: exam=1 round=35 year=2024 period=1 form=A subject=civil_law subject_ko=민법 seq=2 number=2
%% source: Q-Net 감정평가사 기출문제 (cst00309) artlSeq=5213229 file=raw/1008/5213229/1교시 문제지 원본.pdf
%% conversion: kordoc 4.24.51 -> tools/gen1.py (auto); review=proofread-sonnet; flags=none
\begin{question}{2}
신의성실의 원칙에 관한 설명으로 옳지 않은 것은? (다툼이 있으면 판례에 따름)
\begin{choices}
  \item 숙박계약상 숙박업자는 투숙객의 안전을 배려하여야 할 신의칙상 보호의무를 부담한다.
  \item 입원계약상 병원은 입원환자에 대하여 휴대품 도난 방지를 위하여 필요한 적절한 조치를 할 신의칙상 보호의무가 있다.
  \item 기획여행계약상 여행업자는 여행객의 신체나 재산의 안전을 배려할 신의칙상 보호의무를 부담한다.
  \item 계약성립의 기초가 되지 않은 사정의 변경으로 일방당사자가 계약 당시 의도한 계약 목적을 달성할 수 없게 되어 손해를 입은 경우, 그 계약의 효력을 그대로 유지하는 것은 특별한 사정이 없는 한 신의칙에 반한다.
  \item 토지거래허가구역 내의 토지에 관해 허가를 받지 않고 매매계약을 체결한 자가 허가가 없음을 이유로 그 계약의 무효를 주장하는 것은 특별한 사정이 없는 한 신의칙에 반하지 않는다.
\end{choices}
\end{question}
\answer{4}
